% ---
% id: reductions/cor:reference-companion
% kind: cor
% title: "The reference companion of a reduction instance"
% summary: >-
%   The same instance, run on the reference in place of the snap, or on its
%   own snapped blocks where none is handed, is a reduction instance of the
%   emitted shape that squashes that model exactly. Its entries are enclosed
%   from the reference, requested ones where their certificates hold there.
% prominence: main
% status: accepted
% uses: [def:reduction-instance, def:substrate, cor:identity-instance, thm:instrument-instance, cor:register-instance, cor:blockdiag, thm:transfer, thm:chain, def:link, def:chain, def:squashing, thm:licence, cor:entries, def:emitnorm, receiver/lem:enclosure-arithmetic, receiver/def:complex-enclosure, foundations/def:snap, foundations/def:exact-reference, foundations/def:evaluated-off-snap]
% mentions: [thm:emit, lem:row-relaxation, req:reference-companion, reductions/wpfss/cor:flat-companion, reductions/wpfss/rem:published]
% assumes: [A-SNAP]
% refs: []
% provenance: newly-derived
% lean:
% notes: "Part (vi) is stated per writing of an entry in the form eq:instance-entry; the engine forms the companion form of a requested entry at the writing by which it recognised the entry, and a writing whose certificate at the companion is not established gives no entry of the companion."
% ---

An emission is computed on the snapped measurement a declaring instance
states, and the normalizer may move the instance's datum once more
(Theorem~\ref{thm:emit}). Where the declaring instance holds a reference of
the measurement and hands it beside its snap, the same construction can be
run on the reference instead: the instance's datum formed from the reference
blocks, with no normalizer. Where it hands none, the same construction run on
the snapped blocks themselves gives the producer's exact construction, the
\tref{handed datum}, which the normalizer has not touched: the snap of the
emission is then bypassed though the instance's own snap is not. Nothing has
to be proved again for that, because
every part of a \tref{reduction instance} is established at every \tref{admissible datum};
the one thing to supply is that the datum formed from the reference is
admissible, which is hypothesis (H-REF) below, and the substrate in which the
construction is read, since reference entries are in general irrational.
What results is no declaration a consumer solves in its own right
(Requirement~\ref{req:reference-companion}); its entries, which are formed
from the reference, are emitted in binary64 by
Lemma~\ref{lem:row-relaxation}. Part~(iv) discharges (H-REF) for the identity
instance and passes it through branching; the flat instance discharges it in
Corollary~\ref{cor:flat-companion}. Part~(vi) carries to the companion the
entries a caller requests beside the instance's own: their operators are
formed at the companion's datum by the same formula, and their certificates
are re-established there, since a certificate decided on the emitted
numbers says nothing about the reference; for the operator of
\eqref{eq:two-sector-entry} the flat instance does so in
Corollary~\ref{cor:flat-companion}.

\begin{corollary}\label{cor:reference-companion}
Let $\mathcal{J}$ be a \tref{reduction instance}
(Definition~\ref{def:reduction-instance}) whose \tref{handed datum} $A$ is
formed, by the rule of its part (I1), from a family $K=\{K^{(b)}_j\}$ of
measurement blocks, over blocks $b$ of a \tref{sector map} and outcomes $j$ of an
alphabet (one such family per setting where the blocks are given per
setting), that a declaring instance states as a \tref{snap}
(\xref{foundations/def:snap}) of the measurement of a device. Let
$G=\{G^{(b)}_j\}$ be a family over the same blocks, alphabet and dimensions,
every block of which is Hermitian and positive semidefinite and satisfies
exactly every completeness identity that the declaration of $K$ states of
$K$ on that block. In the use this corollary is for, $G^{(b)}$ is the block
of the \tref{handed reference} (\xref{foundations/def:evaluated-off-snap})
where the declaring instance hands one, and $K^{(b)}$ elsewhere, so that
$G=K$ where it hands none. Suppose
\begin{itemize}[leftmargin=4.2em]
\item[(H-REF)] $\mathcal{J}$ is a reduction instance over the substrate
  $\Fld=\mathbb{R}$ (Definition~\ref{def:substrate}), and the rule of its
  part (I1) forms, from every family with the properties just listed for $G$,
  a datum in $\mathcal{D}_{\mathrm{adm}}$.
\end{itemize}
Let $A^{\mathrm{ref}}$ be the datum that rule forms from $G$ in place of
$K$, and call the instance at $A^{\mathrm{ref}}$ --- the same $\Bsys$, (H),
sectors, supports, \tref[scalar value]{scalar values}, constructor and entry rule, with
$A^{\mathrm{ref}}$ in place of the handed datum --- the
\emph{reference companion} of $\mathcal{J}$ at $G$. Then:
\begin{enumerate}[label=(\roman*),leftmargin=3.0em]
\item \emph{(an admissible target over the emitted shape)}
  $A^{\mathrm{ref}}\in\mathcal{D}_{\mathrm{adm}}$, the reference companion
  is a reduction instance, and its target $\{F_v[A^{\mathrm{ref}}]\}_{v\in\Vset}$
  is a POVM on $\Qsys$, sector-block-diagonal by construction. Its alphabet
  with classes and weights, its sector map (the sector list, its order, each
  sector's dimension and kind), its scalar values and its party partition are
  those of the target the instance has at every admissible datum, the
  \tref{declared datum} the normalizer produces included, and it differs
  from such a target on the \tref[data sector]{data sectors} alone, where it
  carries $A^{\mathrm{ref}}$. Where $G=K$, $A^{\mathrm{ref}}=A$: the
  reference companion is $\mathcal{J}$ at its handed datum;
\item \emph{(an exact squash of the companion's model)}
  $\big(M[A^{\mathrm{ref}}];O[A^{\mathrm{ref}}]\big)\dsq_{\Lambda[A^{\mathrm{ref}}]}
  \big(F[A^{\mathrm{ref}}];\hat O[A^{\mathrm{ref}}]\big)$, so that
  $\Lambda[A^{\mathrm{ref}}]^{\dagger}(F_v[A^{\mathrm{ref}}])=M_v[A^{\mathrm{ref}}]$
  exactly for every $v$; every certificate of the entry list
  $\CList[A^{\mathrm{ref}}]$ is discharged, and the \tref{class-marginal identity}
  \eqref{eq:instance-classmarginal} holds at $A^{\mathrm{ref}}$ where the
  instance's flag says it is discharged; and Theorem~\ref{thm:transfer} and
  Theorem~\ref{thm:chain} hold of the one-link \tref{receiver segment}
  (Definition~\ref{def:link}) whose \tref{source family} is
  $(\Bsys,M[A^{\mathrm{ref}}])$, whose target is
  $(\Qsys,F[A^{\mathrm{ref}}])$ and whose \tref{channel witness} is
  $\Lambda[A^{\mathrm{ref}}]$, with $\textsf{status}=\textsf{exact}$;
\item \emph{(the companion is known through enclosures)} let every block of
  $G$ be read with an \tref{entrywise enclosure}
  (\xref{foundations/def:exact-reference}), with radius zero on every block
  where $G^{(b)}=K^{(b)}$. Where the rule of (I1) forms each entry of the
  datum as a finite expression in entries of the blocks and in numbers the
  instance was handed, and every reciprocal node of that expression receives
  an enclosure by part (iv) of \xref{receiver/lem:enclosure-arithmetic}, as
  it does where the expression has none, the rules of that lemma applied to
  that expression give a \tref{complex enclosure} of every element of
  $A^{\mathrm{ref}}$ whose centre is formed by the same expression from the
  centres and whose radius is zero at every entry whose inputs are all
  enclosed with radius zero; part (v) of that lemma, applied element by
  element, makes it an entrywise enclosure of $A^{\mathrm{ref}}$, whose
  centre is the Hermitian part of that centre and whose radius is zero at an
  entry $(i,l)$ where the radii at $(i,l)$ and $(l,i)$ both are. Every entry operator
  $\hat O_k[A^{\mathrm{ref}}]$ of \eqref{eq:instance-entry}, a combination
  of the blocks $\Pi_sF_v[A^{\mathrm{ref}}]\Pi_s$ with real coefficients
  $\alpha_{k,v,s}$ plus a combination $Z_k$ of sector projectors, is then
  enclosed entrywise by the same combination of the centres, with radius
  $\sum_{v}|\alpha_{k,v,s}|\,R^{(s)}_{v,il}$ at the entry $(i,l)$ of each
  data sector $s$, where $(C,R)$ is any entrywise enclosure of the datum and
  $R^{(s)}_{v,il}$ is the radius of its element $v$ at that entry, and
  radius zero on every \tref{scalar sector};
\item \emph{(where (H-REF) holds)} (H-REF) holds for the identity instance of
  Corollary~\ref{cor:identity-instance}; and it holds for the branching
  instance of Theorem~\ref{thm:instrument-instance}, and for the register
  instance of Corollary~\ref{cor:register-instance}, whenever it holds for
  every branch instance whose datum is formed from blocks of $K$, the data of
  the other branches being formed from no block of $K$;
\item \emph{(what the companion rests on, and only there)} no part of
  (i)--(iv) rests on \xref{foundations/ass:snapped-measurement}. Let $B$ be
  the blocks on which $G^{(b)}=K^{(b)}$ that the record of the run names as
  resting on that assumption. Where $B$ is empty this part asserts nothing.
  Where it is not, let the assumption hold: then on every block in $B$ the
  block of the measurement the device performs is the snapped block
  $K^{(b)}$, with the entries recorded for the snap, in place of the block
  of the declaring instance's model, so that the companion is formed there
  from the block the device then measures; and no claim relates a key length
  computed from the companion to that model on the blocks in $B$;
\item \emph{(entries requested beside the instance's own)} let
  $(\hat O,\bowtie,c)$ be an entry emitted at a request that carries entries
  (Corollary~\ref{cor:entries}), written in the form
  \eqref{eq:instance-entry} with coefficients $\alpha_{v,s}$ and a
  combination $Z$ of sector projectors,
  $\hat O=\sum_{v\in\Vset}\sum_{s\in\SecSet}\alpha_{v,s}\Pi_sF_v[A']\Pi_s+Z$
  at the \tref{declared datum} $A'$, and call
  $(\hat O^{\mathrm{ref}},\bowtie,c)$, with
  \begin{equation}\label{eq:companion-form}
    \hat O^{\mathrm{ref}}\;\coloneqq\;\sum_{v\in\Vset}\sum_{s\in\SecSet}
      \alpha_{v,s}\,\Pi_sF_v[A^{\mathrm{ref}}]\Pi_s\;+\;Z ,
  \end{equation}
  its \emph{companion form} at that writing. Part~(iii) encloses
  $\hat O^{\mathrm{ref}}$. Where $\hat O$ is a sum of sector projectors
  (Corollary~\ref{cor:entries}\,(i)), written with every $\alpha_{v,s}=0$,
  $\hat O^{\mathrm{ref}}=\hat O$ and
  $\Lambda[A^{\mathrm{ref}}]^{\dagger}(\hat O)\succeq0$. Where it is the
  class-marginal entry $\big(\sum_{v\in\Vset_t}F_v,=,q_t\big)$ of
  Corollary~\ref{cor:entries}\,(ii), written with $\alpha_{v,s}=1$ for
  $v\in\Vset_t$, $0$ otherwise, and $Z=0$, $\hat O^{\mathrm{ref}}=\hat O$,
  and where the instance discharges \eqref{eq:instance-classmarginal},
  $\Lambda[A^{\mathrm{ref}}]^{\dagger}(\hat O^{\mathrm{ref}})=q_t\,\idon{\Bsys}$
  exactly. For any other entry this part asserts no certificate at
  $A^{\mathrm{ref}}$. In every case, where
  $\Lambda[A^{\mathrm{ref}}]^{\dagger}(\hat O^{\mathrm{ref}})\bowtie
  c\,\idon{\Bsys}$ holds,
  $\big(M[A^{\mathrm{ref}}];O[A^{\mathrm{ref}}],
  \Lambda[A^{\mathrm{ref}}]^{\dagger}(\hat O^{\mathrm{ref}})\big)
  \dsq_{\Lambda[A^{\mathrm{ref}}]}
  \big(F[A^{\mathrm{ref}}];\hat O[A^{\mathrm{ref}}],\hat O^{\mathrm{ref}}\big)$,
  and the companion form is an entry carrying the spectral certificate
  \eqref{eq:certificate} at its constant.
\end{enumerate}
\end{corollary}

\begin{proof}
\emph{Idea.} The \tref{stability clause} of Definition~\ref{def:reduction-instance}
was stated so that the instance at any admissible datum is again an
instance; (H-REF) supplies an admissible datum formed from $G$, and every
conclusion is a part of the instance read at that datum.

(i) The family $G$ has the properties (H-REF) lists, so
$A^{\mathrm{ref}}\in\mathcal{D}_{\mathrm{adm}}$, the set of admissible data
over the substrate $\Fld=\mathbb{R}$, which Definition~\ref{def:substrate}
admits: it leaves the field open between $\mathbb{Q}$ and $\mathbb{R}$, and
entries of a reference may be irrational. The definition's closing sentence
says that the instance at an admissible datum is a reduction instance
whenever $\mathcal{J}$ is, by the stability clause: (H) reads no entry of
the datum, so it holds of $M[A^{\mathrm{ref}}]$ as it holds of the \tref{handed
model}, and (I2)--(I4) are established at every admissible datum. By the
first consequence of Definition~\ref{def:reduction-instance},
$\{F_v[A^{\mathrm{ref}}]\}$ is a POVM on $\Qsys$, its data sectors carrying
$A^{\mathrm{ref}}$ and its scalar sectors the scalar values; it is assigned
one sector block at a time by \eqref{eq:instance-target}, so it is
sector-block-diagonal by construction. The alphabet, the classes and
weights, the sector list with its order, dimensions and kinds, the scalar
values and the declared party label are data of the instance that the
instance at a datum keeps by definition; only the datum changes, and the
datum occupies the data sectors alone. Where $G=K$, the rule of (I1) is
applied to the family it formed the handed datum from, so
$A^{\mathrm{ref}}=A$.

(ii) The relation is (I3) of the reference companion, which (i) shows to be
a reduction instance; its first clause is
$\Lambda[A^{\mathrm{ref}}]^{\dagger}(F_v[A^{\mathrm{ref}}])
=M_v[A^{\mathrm{ref}}]$ by Definition~\ref{def:squashing}. Every
certificate of $\CList[A^{\mathrm{ref}}]$ and, where flagged, the
class-marginal identity are (I4) at $A^{\mathrm{ref}}$, discharged there by
the stability clause. Theorem~\ref{thm:transfer} needs only the relation.
The one-link segment with source family $(\Bsys,M[A^{\mathrm{ref}}])$,
target $(\Qsys,F[A^{\mathrm{ref}}])$, channel witness
$\Lambda[A^{\mathrm{ref}}]$, the transported pair
$(\hat O[A^{\mathrm{ref}}],O[A^{\mathrm{ref}}])$, the \tref{reachable-set statement}
of that channel and the entry list $\CList[A^{\mathrm{ref}}]$ is a
receiver-stage link of Definition~\ref{def:link} and a chain of
Definition~\ref{def:chain} with an empty \tref{source segment}, and
Theorem~\ref{thm:chain}\,(iv) gives it $\textsf{status}=\textsf{exact}$,
the relation carrying no error term. The sector pinching in front of
$\Lambda_0[A^{\mathrm{ref}}]$ changes none of this
(Corollary~\ref{cor:blockdiag}\,(i)).

(iii) Each entry of $A^{\mathrm{ref}}$ is the value of the expression the
rule of (I1) uses, at the entries of $G$ and at the handed numbers. Those
handed numbers that are not entries of a block are the same as for $A$ and
are known exactly; an exactly known number is enclosed by itself with radius
zero (\xref{receiver/def:complex-enclosure}), and an entrywise enclosure of a
block is a complex enclosure of each of its entries. The rules (i)--(iii) of
\xref{receiver/lem:enclosure-arithmetic} enclose sums, products, conjugates
and, where the condition of its part (iii) is met, reciprocals of enclosed
numbers, with centres formed by the same operation on the centres, and keep
an enclosure exact where the inputs are; by its part (iv), applied along the
expression under the hypothesis that every reciprocal node receives an
enclosure, they give a complex enclosure of every entry of
$A^{\mathrm{ref}}$, with the centre and the radii stated. Every element of
$A^{\mathrm{ref}}$ is a Hermitian matrix, the datum being admissible (i), so
part (v) of that lemma turns the complex enclosure of each element into an
entrywise enclosure in the sense of \xref{foundations/def:exact-reference},
with Hermitian centre $\tfrac12(X_c+X_c^{\dagger})$ and symmetric radii
$\tfrac12(r+r^{\transpose})$, exact at an entry $(i,l)$ where
$r_{il}=r_{li}=0$. For an entry
operator, $\hat O_k[A^{\mathrm{ref}}]$ restricted to a data sector $s$ is
$\sum_v\alpha_{k,v,s}A^{\mathrm{ref},(s)}_v$ plus the coefficient of $\Pi_s$
in $Z_k$ times $\idon{d_s}$, exactly known; at each entry the triangle
inequality gives $\big|\sum_v\alpha_{k,v,s}(A^{\mathrm{ref},(s)}_v
-C^{(s)}_v)_{il}\big|\le\sum_v|\alpha_{k,v,s}|\,R^{(s)}_{v,il}$, with
$(C,R)$ any entrywise enclosure of the datum. On a scalar sector the operator is a
combination of the scalar values, which are the instance's and exact.

(iv) For the identity instance the rule of (I1) is the identity: the datum
is the declared family read as a sector-block list, and it is admissible
exactly when each block is Hermitian and positive semidefinite and the class
sums are $q_t$ times the identity, the completeness identities its
declaration states. A family $G$ with the properties listed has them, so the
datum it forms is admissible. That the identity instance is an instance over
$\mathbb{R}$ is the proof of Corollary~\ref{cor:identity-instance}, which
reads of the substrate its field operations and its order and of a datum
only admissibility: its channel, the sector pinching, does not depend on
the datum. For the branching and the register instances,
$\mathcal{D}_{\mathrm{adm}}=\prod_j\mathcal{D}^{(j)}_{\mathrm{adm}}$, a
composite datum being admissible exactly when every branch datum is (Step~1
of the proof of Theorem~\ref{thm:instrument-instance}, which
Corollary~\ref{cor:register-instance} applies); replacing $K$ by $G$ changes
only the branch data formed from blocks of $K$, each admissible by (H-REF)
for its branch, and leaves the others as they are. Both results establish
the composite instance from its branch instances by arguments that read of
the substrate its field operations and its order, so the composite is an
instance over $\mathbb{R}$ when its branches are.

(v) No step of (i)--(iv) reads the assumption. Where $B$ is not empty, each
block in $B$ is a block of the declaring instance's snap of the measurement
of a device, the data $D$ of that snap being what the instance states as
that measurement or a reference of it, so clause (i) of
\xref{foundations/ass:snapped-measurement} applies to it: the block of the
measurement the device performs is the snapped family there, with the
entries recorded for the snap, in place of the model. On such a block
$G^{(b)}=K^{(b)}$, so the companion is formed there from that snapped block.
A nearness of the snapped block and the model relates no key lengths, and
no claim relates the two on a block in $B$, which is a \tref{moved} block of
the snap, as the assumption's $B$ is defined; so no claim relates the
companion's key length to the model there.

(vi) The companion form is an operator of the form
\eqref{eq:instance-entry} at the datum $A^{\mathrm{ref}}$, which
(iii) encloses. Where every $\alpha_{v,s}$ is zero, \eqref{eq:companion-form}
reads no datum, so $\hat O^{\mathrm{ref}}=Z=\hat O$; $Z$ is a sum of sector
projectors, each in the \tref{transported set} of the companion by (I3) at
$A^{\mathrm{ref}}$, and $\Lambda[A^{\mathrm{ref}}]^{\dagger}$ is a positive
map (the third consequence of Definition~\ref{def:reduction-instance}), so
$\Lambda[A^{\mathrm{ref}}]^{\dagger}(Z)\succeq0$, as
Corollary~\ref{cor:entries}\,(i) argues at $A'$. For the class-marginal
entry, $F_v[A]$ is sector-block-diagonal at every admissible datum (i), so
$\sum_s\Pi_sF_v[A]\Pi_s=F_v[A]$ and
$\hat O^{\mathrm{ref}}=\sum_{v\in\Vset_t}F_v[A^{\mathrm{ref}}]$. On a scalar
sector $s$ the block of $\sum_{v\in\Vset_t}F_v$ is formed from the scalar
values, which are the same at both data (i). On a data sector $s$ it is
$\sum_{v\in\Vset_t}A^{(s)}_v$ at the datum it is read at, the elements
\tref{outside the sector support} being zero; an admissible datum is
\tref{class complete}, so this sum is $q_t\idon{d_s}$ at $A^{\mathrm{ref}}$,
admissible by (i), and at $A'$, admissible as an element of
$\mathcal{D}_{\kappa}$ (Definition~\ref{def:emitnorm}). The two operators
therefore agree block by block, and $\hat O^{\mathrm{ref}}=\hat O$. Where
the instance discharges \eqref{eq:instance-classmarginal}, (ii) gives that
identity at $A^{\mathrm{ref}}$ and
$\Lambda[A^{\mathrm{ref}}]^{\dagger}(F_v[A^{\mathrm{ref}}])
=M_v[A^{\mathrm{ref}}]$, so by linearity
$\Lambda[A^{\mathrm{ref}}]^{\dagger}(\hat O^{\mathrm{ref}})
=\sum_{v\in\Vset_t}M_v[A^{\mathrm{ref}}]=q_t\idon{\Bsys}$, which is
\eqref{eq:certificate} with $\bowtie$ the equality at $c=q_t$. For any
other entry the certificate was established at $A'$, for the numbers
emitted there (Corollary~\ref{cor:entries}\,(iii)), and a certificate of an operator
formed from one datum is no statement about the operator formed from
another; so none is asserted. A writing of $\hat O$ in the form
\eqref{eq:instance-entry} need not be unique, and two writings may give
different companion forms; each is an entry wherever its own certificate
holds, which is all this part claims of it. Finally, where
$\Lambda[A^{\mathrm{ref}}]^{\dagger}(\hat O^{\mathrm{ref}})\bowtie
c\,\idon{\Bsys}$ holds, (ii) is a \tref{squashing relation} at
$A^{\mathrm{ref}}$, and adjoining the pair
$\big(\Lambda[A^{\mathrm{ref}}]^{\dagger}(\hat O^{\mathrm{ref}}),
\hat O^{\mathrm{ref}}\big)$ to its \tref{carried operators} keeps
\eqref{eq:squashing}, whose clause for that pair holds by the choice of its
physical operator (Definition~\ref{def:squashing}); the certificate is then
\eqref{eq:certificate} for the adjoined carried operator.
\end{proof}
